\section{Literal relative normal associated graded of an original ideal}
Let $B$ be the entire tangential holomorphic germ ring, $R$ the original
ambient holomorphic germ ring and $J$ the actual ideal of the normal
coordinates. Fix an arbitrary original whole ideal $I\subseteq R$.
No caller kernel, finite relation, flattening, basis or generic-freeness
certificate is an input to the following constructions.

\begin{theorem}[Actual original normal monomials and literal graded layers]
\label{hir:relative-normal-graded-literal-layers}
\lean{Hironaka.HironakaRelativeNormalAssociatedGraded.normalExponentDegree_eq_degree,Hironaka.HironakaRelativeNormalAssociatedGraded.normalExponentMonomial_mem_pow,Hironaka.HironakaRelativeNormalAssociatedGraded.originalNormalGradedPiece,Hironaka.HironakaRelativeNormalAssociatedGraded.originalNormalGradedPieceAddCommGroup,Hironaka.HironakaRelativeNormalAssociatedGraded.originalNormalAssociatedGraded,Hironaka.HironakaRelativeNormalAssociatedGraded.originalNormalPowerImage,Hironaka.HironakaRelativeNormalAssociatedGraded.originalNormalPowerImage_succ_le,Hironaka.HironakaRelativeNormalAssociatedGraded.originalNormalGradedStepMap,Hironaka.HironakaRelativeNormalAssociatedGraded.originalNormalGradedStepMap_surjective,Hironaka.HironakaRelativeNormalAssociatedGraded.ker_originalNormalGradedStepMap,Hironaka.HironakaRelativeNormalAssociatedGraded.originalNormalGradedPieceQuotientEquiv}
\leanok
\uses{hir:relative-normal-basis-actual-entire-quotient-basis,hir:relative-normal-ideal-actual-coordinates}
The actual normal monomial $z^\alpha$ belongs to $J^{|\alpha|}$.
The literal degree-$q$ layer is the entire $B$-submodule
$L_q=\operatorname{image}(J^q\longrightarrow R/(I+J^{q+1}))$.
It is precisely the quotient of the image of $J^q$ in $R/I$ by the image
of $J^{q+1}$, with the true original representative dictionary. The
associated graded module is the actual direct sum $\bigoplus_q L_q$.
\end{theorem}
\begin{proof}
Each normal coordinate belongs to its genuine generated ideal; its $k$th
power belongs to the $k$th ideal power. Multiply these actual memberships
and use the finite product identity $\prod_iJ^{\alpha_i}=J^{\sum_i\alpha_i}$.
Use the literal quotient map and submodule image to define $L_q$. Its
representatives are precisely original germs in $J^q$. Two such germs
have the same layer class exactly when their difference lies in
$I+J^{q+1}$. Passing first to $R/I$ gives the same dictionary, since a
representative in $J^q$ maps to the next power image exactly when its
difference from some actual $J^{q+1}$ germ lies in $I$. The first
isomorphism theorem gives the actual layer identification and hence the
literal direct sum. This includes $q=0$ and zero/unit original ideals.
\end{proof}

\begin{theorem}[Genuine homogeneous normal polynomial surjection]
\label{hir:relative-normal-graded-homogeneous-surjection}
\lean{Hironaka.HironakaRelativeNormalAssociatedGraded.normalPolynomialEval_mem_pow_of_homogeneous,Hironaka.HironakaRelativeNormalAssociatedGraded.exists_homogeneous_normal_polynomial_representative,Hironaka.HironakaRelativeNormalAssociatedGraded.homogeneousNormalSymbolMap,Hironaka.HironakaRelativeNormalAssociatedGraded.homogeneousNormalSymbolMap_surjective,Hironaka.HironakaRelativeNormalAssociatedGraded.normalAssociatedSymbolMap,Hironaka.HironakaRelativeNormalAssociatedGraded.normalAssociatedSymbolMap_surjective,Hironaka.HironakaRelativeNormalAssociatedGraded.normalAssociatedSymbolMap_component,Hironaka.HironakaRelativeNormalAssociatedGraded.normalAssociatedSymbolMap_eq_zero_iff}
\leanok
\uses{hir:relative-normal-graded-literal-layers,hir:relative-normal-ideal-actual-coordinates,hir:relative-normal-basis-actual-coefficient-separation}
Evaluate an actual degree-$q$ homogeneous normal polynomial in the
original normal coordinates and then take its actual class in $L_q$.
This is a true $B$-linear surjection onto $L_q$. Summing these maps after
the genuine polynomial homogeneous decomposition gives a true surjection
$B[Z_1,\ldots,Z_t]\twoheadrightarrow\bigoplus_qL_q$.
Its whole kernel is exactly the condition that the evaluated $q$th
homogeneous part belongs to $I+J^{q+1}$ for every $q$.
\end{theorem}
\begin{proof}
Homogeneous polynomial support has degree $q$, so actual monomial
membership proves that its entire evaluation belongs to $J^q$.
Given an original representative $g\in J^q$, use the genuine distinct
normal-monomial representative modulo $J^{q+1}$. Its evaluated polynomial
also lies in $J^q$, because the actual remainder is in $J^{q+1}\subseteq
J^q$. Whole polynomial coefficient separation forces every coefficient
of degree less than $q$ to vanish as an entire base germ. All its
remaining degrees are exactly $q$, giving the required homogeneous
representative. The true finite homogeneous decomposition of a polynomial
and the actual universal property of the direct sum construct the full
map. Surjectivity on each literal layer implies surjectivity on the true
direct sum. Equality to zero in a layer is precisely whole membership in
$I+J^{q+1}$, proving the asserted complete kernel condition.
\end{proof}

\begin{theorem}[True original initial symbols and the complete homogeneous kernel ideal]
\label{hir:relative-normal-graded-whole-initial-ideal}
\lean{Hironaka.HironakaRelativeNormalAssociatedGraded.normalAssociatedSymbolMap_eq_zero_of_homogeneous,Hironaka.HironakaRelativeNormalAssociatedGraded.homogeneous_normalPolynomialEval_mul_mem_remainder,Hironaka.HironakaRelativeNormalAssociatedGraded.normalAssociatedSymbolMap_mul_eq_zero,Hironaka.HironakaRelativeNormalAssociatedGraded.originalNormalInitialIdeal,Hironaka.HironakaRelativeNormalAssociatedGraded.originalNormalInitialIdeal_eq_ker,Hironaka.HironakaRelativeNormalAssociatedGraded.homogeneous_normal_polynomial_unique,Hironaka.HironakaRelativeNormalAssociatedGraded.originalNormalInitialSymbol,Hironaka.HironakaRelativeNormalAssociatedGraded.originalNormalInitialSymbol_homogeneous,Hironaka.HironakaRelativeNormalAssociatedGraded.originalNormalInitialSymbol_remainder,Hironaka.HironakaRelativeNormalAssociatedGraded.originalNormalRelationIndex,Hironaka.HironakaRelativeNormalAssociatedGraded.originalNormalRelation,Hironaka.HironakaRelativeNormalAssociatedGraded.originalNormalRelation_mem_initialIdeal,Hironaka.HironakaRelativeNormalAssociatedGraded.exists_original_germ_homogeneous_kernel_relation,Hironaka.HironakaRelativeNormalAssociatedGraded.originalNormalInitialIdeal_eq_span_original_relations,Hironaka.HironakaRelativeNormalAssociatedGraded.originalNormalInitialIdeal_homogeneous}
\leanok
\uses{hir:relative-normal-graded-homogeneous-surjection}
For every actual $g\in I\cap J^q$, its true degree-$q$ normal initial
symbol is the $q$th homogeneous part of its actual polynomial
representative modulo $J^{q+1}$. Let $H$ be the ideal generated by all
these actual homogeneous symbols. Then $H$ is homogeneous and is the
COMPLETE whole kernel of the preceding polynomial surjection.
\end{theorem}
\begin{proof}
The initial symbol is homogeneous by the actual polynomial projection,
and its evaluation differs from $g$ by a true element of $J^{q+1}$.
It therefore maps to zero. Multiplication of homogeneous polynomials
corresponds to multiplication of their original evaluated germs: an
actual error in $I+J^{q+1}$ times a degree-$r$ representative in $J^r$
lies in $I+J^{q+r+1}$. Decompose every multiplier into its actual finite
homogeneous parts to obtain closure of the complete kernel under arbitrary
polynomial multiplication. Addition and entire base scalar multiplication
are actual linear operations. This proves an actual ideal, and stability
under all homogeneous projections proves homogeneity. Conversely, a
homogeneous $P$ of degree $q$ in the kernel has actual evaluation
$g+h$ with $g\in I$ and $h\in J^{q+1}$. Since its evaluation and $h$
belong to $J^q$, the entire original $g$ lies in $I\cap J^q$. Whole
coefficient separation modulo $J^{q+1}$ identifies $P$ with the actual
initial symbol of $g$. The finite homogeneous decomposition then expresses
every complete kernel element in the ideal generated by these actual
symbols. Both inclusions are proved with whole germs and actual original
ideal memberships.
\end{proof}

\begin{theorem}[Internally finite actual homogeneous original relations]
\label{hir:relative-normal-graded-internal-finite-relations}
\lean{Hironaka.HironakaRelativeNormalAssociatedGraded.exists_finite_originalNormalRelations,Hironaka.HironakaRelativeNormalAssociatedGraded.originalNormalAssociatedGradedQuotientEquiv,Hironaka.HironakaRelativeNormalAssociatedGraded.exists_finite_originalNormalAssociatedGraded_presentation}
\leanok
\uses{hir:relative-normal-graded-whole-initial-ideal,hir:relative-normal-ideal-actual-original-matrices}
An actual finite family of original germs $g_i\in I\cap J^{q_i}$
internally supplies homogeneous relations generating the complete kernel
$H$. Thus the literal original associated graded module has a true finite
homogeneous polynomial presentation, valid in all degrees simultaneously.
\end{theorem}
\begin{proof}
The genuine entire base holomorphic germ ring is Noetherian. The true
polynomial ring in the finite actual normal coordinate family is therefore
Noetherian by Hilbert's basis theorem. Its actual ideal $H$ is finitely
generated. Apply finite extraction to its proved span by all actual
original initial symbols. This chooses a finite SUBSET of that actual
symbol family, rather than replacing it by caller relations or arbitrary
nonhomogeneous polynomial generators. Choose the original germ and degree
witness for each of these finitely many actual symbols. The complete
kernel equality then identifies the polynomial quotient by these genuine
finite homogeneous relations with the literal original direct sum of
layers. No generic-freeness or flatness result is inferred from finite
presentation alone.
\end{proof}

\begin{remark}
A genuine nonzero common denominator making this entire graded module free,
and a true uniform analytic neighbourhood realizing the required original
stalks, remain subsequent producer obligations. Separate degree-wise
localizations and finite presentation do not settle all-degree normal
flatness. Higher HS descent and general G8 remain OPEN.
\end{remark}
